\begin{table}[htbp]
\centering
\caption{Largest member-party contributions to the within- and between-district components}
\label{tab:coalition-party-component-extremes}
\footnotesize
\setlength{\tabcolsep}{3.5pt}
\renewcommand{\arraystretch}{1.12}
\begin{tabularx}{\textwidth}{@{}L{0.25\textwidth}*{4}{>{\raggedright\arraybackslash}X}@{}}
\toprule
Case & \multicolumn{2}{c}{Within-district \(A_i\)} & \multicolumn{2}{c}{Between-district \(B_i\)} \\
\cmidrule(lr){2-3}\cmidrule(l){4-5}
 & Positive & Negative & Positive & Negative \\
\midrule
\multicolumn{5}{@{}l}{\textit{Observed cabinet inversions}} \\ \addlinespace[2pt]
2014/14-05 & \mbox{PR} (5.04)\newline \mbox{PSD} (4.65) & \mbox{PT} (-1.22)\newline \mbox{PCdoB} (-0.57) & \mbox{PMDB} (4.22)\newline \mbox{PDT} (1.56) & \mbox{PT} (-1.21)\newline \mbox{PR} (-0.74) \\
2022/22-01 & \mbox{UNIÃO} (9.36)\newline \mbox{PT} (9.07) & \mbox{PSB} (-5.55)\newline \mbox{PSOL} (-3.40) & \mbox{UNIÃO} (1.75)\newline \mbox{PDT} (1.02) & \mbox{PSOL} (-2.66)\newline \mbox{PT} (-2.13) \\
\addlinespace
\multicolumn{5}{@{}l}{\textit{Minimal exact-connected ideological inversions (\(k=0\))}} \\ \addlinespace[2pt]
2014/PSB--PTN & \mbox{PSD} (4.65)\newline \mbox{PTB} (4.08) & \mbox{PT DO B} (-3.61)\newline \mbox{PV} (-1.40) & \mbox{PMDB} (4.22)\newline \mbox{PMN} (0.36) & \mbox{PSDB} (-5.01)\newline \mbox{PV} (-1.16) \\
2014/PTB--PR & \mbox{PR} (5.04)\newline \mbox{PSD} (4.65) & \mbox{PT DO B} (-3.61)\newline \mbox{PRTB} (-1.91) & \mbox{PMDB} (4.22)\newline \mbox{PROS} (0.76) & \mbox{PSDB} (-5.01)\newline \mbox{PR} (-0.74) \\
2014/PT DO B--PSDC & \mbox{PR} (5.04)\newline \mbox{PSD} (4.65) & \mbox{PT DO B} (-3.61)\newline \mbox{PRTB} (-1.91) & \mbox{PMDB} (4.22)\newline \mbox{PROS} (0.76) & \mbox{PSDB} (-5.01)\newline \mbox{PRB} (-2.17) \\
2014/SOLIDARIEDADE--PSL & \mbox{PR} (5.04)\newline \mbox{PSD} (4.65) & \mbox{PSL} (-3.30)\newline \mbox{PRTB} (-1.91) & \mbox{PMDB} (4.22)\newline \mbox{PROS} (0.76) & \mbox{PSDB} (-5.01)\newline \mbox{PRB} (-2.17) \\
2022/MDB--UNIÃO & \mbox{UNIÃO} (9.36)\newline \mbox{MDB} (5.12) & \mbox{PTB} (-5.94)\newline \mbox{PODE} (-5.66) & \mbox{PP} (3.17)\newline \mbox{REPUBLICANOS} (2.94) & \mbox{PSDB} (-0.54)\newline \mbox{MDB} (-0.40) \\
2022/PP--PL & \mbox{PL} (19.07)\newline \mbox{UNIÃO} (9.36) & \mbox{PSC} (-3.71)\newline \mbox{PATRIOTA} (-3.03) & \mbox{PP} (3.17)\newline \mbox{REPUBLICANOS} (2.94) & \mbox{PL} (-5.29)\newline \mbox{NOVO} (-0.55) \\
\bottomrule
\end{tabularx}
\begin{minipage}{0.98\linewidth}
\vspace{0.35em}
\footnotesize\textit{Notes:} Entries show the two largest positive and two most negative contributions in seats; values are rounded to two decimals. Cases are identified by election year and cabinet party-set label or interval endpoints. \(A_i\) and \(B_i\) are ex post accounting contributions: \(d_i=A_i+B_i\), \(A_C=\sum_{i\in C}A_i\), and \(B_C=\sum_{i\in C}B_i\). For 2014 and 2018, party-level attribution is ex post because seats were often allocated through joint electoral lists. These values do not identify party-specific causal effects of electoral rules.

\par{\footnotesize V5 primary sets include 100 flagged provisional days; the underlying UNKNOWN affiliations remain unresolved.}\par{\footnotesize Recorded concrete date and affiliation sensitivities are reported separately.}
\end{minipage}
\end{table}
